\documentclass[a4paper,11pt]{article}
\usepackage[margin=2.5cm]{geometry}
\begin{document}

\begin{center}
{\large\bfseries Relatório de Verificação --- decoder\_scancode\_ascii}\\[4pt]
Lucas Damo \quad$\cdot$\quad Antônio dos Santos\\[4pt]
CUV: \texttt{decoder\_scancode\_ascii\_map.vhd} \quad vs \quad Golden: \texttt{decoder\_scancode\_ascii.sv}
\end{center}

\noindent
Foram aplicados os 36 estímulos alfanuméricos (26 letras + 10 dígitos).
Resultado: 31/36 corretos (letras: 22/26, dígitos: 9/10).
As divergências funcionais encontradas entre o CUV e o Modelo de Referência
estão listadas abaixo.

\vspace{8pt}
\begin{center}
\begin{tabular}{|c|c|c|c|}
\hline
\textbf{Entrada (HEX)} & \textbf{Entrada (DEC)} & \textbf{Resultado} & \textbf{Resultado Obtido} \\
 & & \textbf{Esperado (HEX)} & \textbf{(HEX)} \\
\hline
21 & 33 & 43 (`C') & --- $^{(1)}$ \\ \hline
2B & 43 & 46 (`F') & 66 (`f') \\ \hline
43 & 67 & 49 (`I') & --- $^{(1)}$ \\ \hline
31 & 49 & 4E (`N') & B1 (`$\pm$') \\ \hline
3D & 61 & 37 (`7') & --- $^{(1)}$ \\ \hline
\end{tabular}
\end{center}

\vspace{4pt}
\noindent
$^{(1)}$ O testbench imprime a saída com \texttt{\%s}; nestes casos o caractere obtido não
é imprimível, logo o valor em hexadecimal não pode ser recuperado do log.
Alterar o \texttt{\$display} para \texttt{\%h} para registrar o valor exato.

\end{document}
